% Image credits for the photographs used in this volume (Book 1).
% Machine-readable license records live in images/book1/CREDITS.md; keep
% the two files in sync when adding or removing a photograph.
\chapter*{Crédits des images}

Les dessins de ce livre sont originaux. Les photographies sont
utilisées sous leurs licences libres respectives, avec nos
remerciements à leurs auteurs~:

\begin{itemize}
  \item \emph{John Dalton} (chapitre sur les atomes et les molécules)~:
        gravure de Charles Turner d'après Joseph Allen, 1834, Library of
        Congress via Wikimedia Commons, domaine public.
  \item \emph{La table des éléments de Dalton, 1808} (chapitre sur les
        atomes et les molécules)~: planche de l'ouvrage de John Dalton
        \emph{A New System of Chemical Philosophy}, numérisation de
        l'Internet Archive via Wikimedia Commons, domaine public.
  \item \emph{Bauxite} (chapitre sur les matières premières)~:
        photographie de Werner Schellmann via Wikimedia Commons,
        CC~BY-SA~2.5.
  \item \emph{Granite} (chapitre sur les corps purs)~: photographie
        d'Eurico Zimbres via Wikimedia Commons, CC~BY-SA~2.0~br.
  \item \emph{Sulfate de cuivre anhydre} (chapitre sur l'identification
        des espèces)~: photographie de W.~Oelen via Wikimedia Commons,
        CC~BY-SA~3.0.
  \item \emph{Cristaux de sulfate de cuivre} (chapitre sur
        l'identification des espèces)~: photographie de Stephanb via
        Wikimedia Commons, CC~BY-SA~3.0.
  \item \emph{Antoine-Laurent Lavoisier et Marie-Anne Paulze Lavoisier}
        (chapitre sur les équations ajustées)~: tableau de Jacques-Louis
        David, 1788, Metropolitan Museum of Art via Wikimedia Commons,
        domaine public.
  \item \emph{Téléphone en bakélite} (chapitre sur les matières
        plastiques)~: photographie de Holger Ellgaard via Wikimedia
        Commons, CC~BY-SA~3.0.
  \item \emph{Ernest Rutherford} (chapitre sur l'intérieur de
        l'atome)~: photographie, George Grantham Bain Collection, Library
        of Congress, via Wikimedia Commons, domaine public.
  \item \emph{Halite} (chapitre sur les ions)~: photographie de Robert
        M.~Lavinsky via Wikimedia Commons, CC~BY-SA~3.0.
  \item \emph{La statue de la Liberté} (chapitre sur les métaux et la
        corrosion)~: photographie d'Elcobbola via Wikimedia Commons,
        domaine public.
  \item \emph{Dmitri Mendeleïev} et \emph{son tableau de 1869} (chapitre
        sur le tableau périodique)~: via Wikimedia Commons, domaine
        public.
  \item \emph{La nébuleuse du Crabe} (chapitre sur les éléments)~: NASA,
        ESA, J.~Hester et A.~Loll, domaine public.
  \item \emph{Sodium} (chapitre sur les couches électroniques)~:
        photographie de Dnn87 via Wikimedia Commons, CC~BY-SA~3.0.
  \item \emph{Dichlore dans une ampoule} (chapitre sur les couches
        électroniques)~: photographie de W.~Oelen via Wikimedia Commons,
        CC~BY-SA~3.0.
  \item \emph{Amedeo Avogadro} (chapitre sur la mole)~: lithographie
        d'après un dessin de C.~Sentier, 1856, collection Edgar Fahs
        Smith, via Wikimedia Commons, domaine public.
  \item \emph{Cristaux de neige} (chapitre sur la polarité)~:
        photographies de Wilson Bentley, 1902, via Wikimedia Commons,
        domaine public.
  \item \emph{Louis Pasteur} (chapitre sur la stéréochimie)~:
        photographie de Paul Nadar, via Wikimedia Commons, domaine
        public.
  \item \emph{Cristaux de tartrate d'ammonium} (chapitre sur la
        stéréochimie)~: photographie de Marie-Lan Taÿ Pamart via
        Wikimedia Commons, CC~BY~4.0.
  \item \emph{Svante Arrhenius} (chapitre sur les acides et les
        bases)~: photogravure, 1909, via Wikimedia Commons, domaine
        public.
  \item \emph{Une pile de Volta} (chapitre sur les piles)~: photographie
        de GuidoB via Wikimedia Commons, CC~BY-SA~3.0.
  \item \emph{Wallace Carothers dans son laboratoire} (chapitre sur les
        polymères)~: photographe inconnu, via Wikimedia Commons, domaine
        public.
\end{itemize}

Les liens vers les sources complètes et les adresses des licences de
chaque photographie sont donnés dans
\texttt{images/book1/CREDITS.md}, dans le dépôt du livre. Les
pictogrammes de danger sont ceux du Système général harmonisé de
classification et d'étiquetage des produits chimiques, tels qu'ils sont
dessinés pour le paquet \texttt{ghsystem} de \TeX\ Live.

\bigskip
À côté des photographies et des dessins originaux, certaines figures
sont des illustrations produites par intelligence artificielle, réalisées
avec la génération d'images d'OpenAI à partir de consignes rédigées par
les éditeurs du livre, et vérifiées sur le plan chimique avant d'être
retenues. La consigne complète de chaque image générée est consignée
dans \texttt{images/book1/ai/PROMPTS.md}, dans le dépôt.
\clearpage
